% Mechanical relocation generated from the preservation ledger.
% Existing prose is included byte-for-byte from preserved blocks.

\subsection{A Loop-Oriented Description of a RAPTOR Query}
\label{sec:raptor-loop-description}

A RAPTOR query can be understood as an outer loop over the number of vehicle
boardings, containing a timetable scan and a walking-relaxation loop.  Its
state consists of labels attached to physical stops rather than paths through
a global time-expanded graph.  A label records the earliest retained arrival
at a stop in a given boarding round, together with the journey legs and the
additive attributes needed later, such as waiting, walking, and in-vehicle
time.  The query proceeds as follows.

\begin{enumerate}
  \item \emph{Initialize the access round.}  At the requested departure time,
        the origin receives its initial label and every other stop is initially
        unreachable.  A loop follows the admissible footpaths out of each
        newly reached stop.  Whenever walking produces an earlier label, that
        stop is placed back in the walking queue so that multi-link walking
        connections are closed.  This is round zero because no vehicle has yet
        been boarded.

  \item \emph{Loop over boarding rounds.}  If at most $K$ transfers are
        permitted, the algorithm performs $K+1$ transit rounds.  Round $r$
        starts only from labels retained in round $r-1$ and adds exactly one
        vehicle boarding.  This makes the transfer bound part of the search
        state: a journey requiring another vehicle cannot leak into a round in
        which it is not allowed.

  \item \emph{Loop over scheduled trips and their stops.}  For each trip, its
        stop times are scanned once in travel order.  Before boarding, the scan
        compares the trip departure time at the current stop with the labels
        from the preceding round.  If the passenger can be present in time,
        the trip becomes boardable.  The scan then carries that boarded state
        forward along the same trip.  At every subsequent stop it constructs
        a candidate arrival label, records the boarding and alighting events,
        and retains the better candidate according to the deterministic
        arrival and feature ordering.  If a later stop offers a preferable
        feasible boarding state for the same trip, the carried state is
        replaced before the scan continues.

  \item \emph{Loop over walking connections again.}  Once all trips have been
        scanned for the round, the walking-closure loop starts from the stops
        improved by transit.  Walking changes arrival and walking time but
        consumes no additional boarding, so these labels remain in the same
        round.  They become the possible boarding points for the next round.

  \item \emph{Collect and filter the results.}  After the last permitted
        round, the algorithm loops over each destination and compares its
        labels across rounds.  A label is discarded if another arrives no
        later with no more boardings.  Waiting-time and total-journey-time
        limits are then applied.  The surviving labels form the bounded
        timetable-feasible alternatives, and each retained state supplies its
        vehicle legs and boarding/alighting events.  If no label
        survives, a separate topology-only check distinguishes a disconnected
        pair, a pair exceeding the transfer bound, and a pair for which no
        scheduled journey is feasible at the queried time.
\end{enumerate}

Textbook RAPTOR commonly uses a further \emph{marking} optimization: in the
next round it scans only route patterns serving stops whose labels improved in
the preceding round.  The current implementation preserves the same bounded
round semantics but uses the simpler deterministic variant of scanning every
scheduled trip in every round.  This distinction affects preprocessing time,
not the meaning of the retained journey labels.  In either variant, observed
counts play no role in these loops.  Counts enter only after the retained
journeys have been converted into the fixed measurement-response matrix.

\subsection{A RAPTOR Example}
\label{sec:raptor-example}

RAPTOR is a round-based timetable-routing method.  A round counts vehicle
boardings, not transfers: round zero contains access and walking labels, round
one contains journeys using one vehicle, and round two permits two vehicles
and hence one transfer.  The following small timetable illustrates the logic.

Consider four physical stops $A,B,C,D$ and three scheduled trips:
\begin{center}
\begin{tabular}{llll}
\hline
Trip & Line & Stop sequence with times & Role \\
\hline
$r_1$ & Red   & $A$ 08:00, $B$ 08:10, $C$ 08:20 & first leg \\
$b_1$ & Blue  & $B$ 08:13, $D$ 08:28            & connection \\
$g_1$ & Green & $A$ 08:07, $D$ 08:35            & direct trip \\
\hline
\end{tabular}
\end{center}
Suppose the passenger is ready at $A$ at 07:58 and that the three-minute
connection at $B$ satisfies the minimum-transfer rule.  Initially, round zero
labels $A$ at 07:58.  In round one, scanning trips boardable from $A$ produces
labels $B=08{:}10$ and $C=08{:}20$ through the Red line and $D=08{:}35$
through the direct Green line.  In round two, the label at $B$ makes the Blue
trip boardable at 08:13; it improves the destination label to $D=08{:}28$.
The resulting nondominated alternatives include
\[
  A\xrightarrow{g_1}D
  \quad (0\text{ transfers},\ 08{:}35),
  \qquad
  A\xrightarrow{r_1}B\xrightarrow{b_1}D
  \quad (1\text{ transfer},\ 08{:}28).
\]
Neither alternative dominates the other in both arrival time and number of
transfers.  A zero-transfer limit retains only the direct trip; a one-transfer
limit also admits the faster connection.

The example also shows why routing is time dependent.  A passenger becoming
ready at 08:05 misses $r_1$ but can still board $g_1$, so the connecting
alternative disappears even though the physical OD pair is unchanged.  A
later query could encounter an entirely different set of trips.  RAPTOR
therefore returns timetable-feasible journey summaries for a specified origin
and departure time; it does not estimate OD demand and it does not impose one
all-day route share.

When this search is used to construct a reduced response, each retained
alternative contributes boarding and alighting events for all of its vehicle
legs.  The connecting journey above contributes boardings at $A$ and $B$ and
alightings at $B$ and $D$, whereas the direct journey contributes only a
boarding at $A$ and an alighting at $D$.  This distinction is what allows
event-level counts, when sufficiently informative, to discriminate between
journey alternatives.

\subsection{Why the Reduced Approach Changes the Computational Scale}
\label{sec:why-raptor-scales}

The large difference between detailed assignment and the reduced approach does
not arise from replacing one shortest-path routine by a universally faster
one. The two calculations represent the network differently, are invoked at
different stages, and expose optimization problems of very different
dimensions.

The central computational object is the assignment response, not RAPTOR
itself. Let $A$ map OD-time demand to flows on the links or events of the
detailed assignment, and let $M$ aggregate those flows into the modeled count
rows. With fixed routing,
\[
 \bm x=A\bm f,
 \qquad
 \bm\mu=M\bm x=MA\bm f=B\bm f.
\]
Thus the measurement-response matrix is $B=MA$. It can in principle be built
from either the time-expanded graph or RAPTOR-derived journeys. Once $B$ is
available, estimation no longer needs to know how its columns were obtained.

This also means that a redesigned time-expanded workflow could obtain the same
computational benefit. If it enumerated or summarized a bounded journey set,
projected those journeys directly onto measurements, stored the resulting
compact $B$, and then removed graph traversal from the optimizer, its
estimation stage would be conceptually the same as the reduced workflow. The
present difference is therefore a difference in response construction and
representation, not a mathematical capability exclusive to RAPTOR.

The detailed model constructs a global time-expanded graph. Every scheduled
arrival and departure is represented by an event node, and waiting, riding,
transfer, access, and egress possibilities are represented by links. A forward
assignment propagates passenger flow over this graph. Estimation additionally
requires derivatives of measurements with respect to demand or behavioral
parameters. A direct differentiated construction may therefore retain state
whose scale is governed not only by the number of network events, but also by
the number of destination groups or measurement responses. Repeating these
graph traversals and derivative calculations at many optimization iterations
is what creates the hours-long runtime and the prohibitive memory requirement.

The earlier sharded algorithm already exploited the fact that fixed routing is
linear, but it represented $B$ only \emph{implicitly}. For shard $s$, an
objective evaluation computed
\[
 B_s\bm f_s=M A_s\bm f_s
\]
by loading the corresponding routing state and propagating flows through the
complete time-expanded graph. Its adjoint similarly applied
$A_s^{\mathsf T}M^{\mathsf T}$ by a reverse graph traversal. Routing
probabilities could be reused, but the measurement-response columns were not
constructed and retained. Sharding limited the amount of live state and made
the calculation possible in memory; it did not eliminate the expensive graph
work from each operator product. This is why reducing a run to fewer shards
reduced memory much more clearly than elapsed time.

For small cases, $B=MA$ can be constructed from the detailed graph. At full-
network scale, however, a dense $B$ has one row per measurement and one column
per free OD-time cell and is far too large. Constructing it column by column
would require an enormous number of unit-demand assignments, while a reverse
construction can require event-level derivative state for a very large
collection of measurements. Fixed routing guarantees linearity, but it does
not by itself make explicit matrix construction affordable.

RAPTOR uses a different state representation. It scans timetable routes in a
bounded sequence of rounds, with one additional vehicle boarding per round,
and retains compact labels at stops rather than constructing a global graph of
all timetable events. For one origin and departure time, its working state is
principally determined by the stops, timetable records scanned, and permitted
boarding rounds. It does not carry a measurement derivative for every event
while searching for feasible journeys.

More importantly, the RAPTOR journey representation constructs the relevant
entries of $B$ directly. If journey $j$ for OD-time cell $c$ has fixed share
$p_{j\mid c}$ and contributes $h_{jm}$ events to measurement $m$, then
\[
 B_{mc}=\sum_{j\in\mathcal J_c}p_{j\mid c}h_{jm}.
\]
Unobserved time-expanded links never appear in this construction. Each retained
journey touches only a small number of boarding and alighting rows, so the
response can be stored sparsely and identical columns can be compressed. This
is the practical difference: the graph-based shards repeatedly \emph{apply} an
implicit $B$, whereas the reduced approach builds a compact measurement-level
$B$ once and reuses it.

The two constructions produce the same matrix only when they use equivalent
OD-time definitions, physical-stop mappings, access and transfer rules,
feasible journey sets, generalized costs, route-choice probabilities, and
measurement aggregation. Bounded transfer depth, journey pruning, or different
share rules may instead give
\[
 B_{\mathrm{RAPTOR}}\ne B_{\mathrm{graph}}.
\]
The reduced matrix is exact relative to its retained journeys and fixed shares,
but its agreement with the detailed response is an empirical validation
question. Counts are used to estimate demand after either matrix is fixed; they
do not choose the RAPTOR journeys.

The decisive saving comes from using this search as preprocessing rather than
inside every objective evaluation. Feasible journeys and their attributes are
computed once, converted into fixed expected measurement responses, and then
reused. Estimation operates on those responses. With the conditional gravity
model, it also optimizes a small collection of shared behavioral and production
coefficients instead of one variable per feasible OD-time cell. The resulting
workflow is therefore
\[
 \begin{gathered}
   \text{timetable search once}
   \longrightarrow \text{compact fixed response}\\
   \Downarrow\\[-1mm]
   \text{many inexpensive objective evaluations}
 \end{gathered}
\]
whereas repeated detailed estimation places the expensive network calculation
inside the optimization loop.

Three changes are consequently responsible for laptop-scale estimation:
\begin{enumerate}
  \item explicit construction and reuse of a compact measurement-response
        matrix, rather than repeated application of an implicit graph-based
        operator;
  \item timetable-native, bounded-round routing that makes this compact
        response practical to construct; and
  \item a low-dimensional gravity parameterization instead of an unrestricted
        OD-time parameter vector.
\end{enumerate}
These effects must not be attributed to RAPTOR alone. Searching every origin,
departure instant, and destination can still require substantial preprocessing,
and storing uncompressed journey or measurement responses can still be large.
The approach remains practical because routing is bounded, preprocessing is
reused, responses are sparse or compressed, and detailed assignment is reserved
for post-fit validation.
